\documentclass[10pt,a4paper]{amsart}
\usepackage{amsmath,amssymb,amsthm,graphicx}
\usepackage{verbatim}
\newenvironment{reference}{\comment}{\endcomment}
%\newenvironment{reference}{}{}
\newenvironment{slogan}{\comment}{\endcomment}
\newenvironment{history}{\comment}{\endcomment}

% For commutative diagrams we use Xy-pic
\usepackage[all]{xy}

% We use 2cell for 2-commutative diagrams.
\xyoption{2cell}
% Theorem environments.
%
\theoremstyle{plain}
\newtheorem{theorem}[subsection]{Theorem}
\newtheorem{proposition}[subsection]{Proposition}
\newtheorem{lemma}[subsection]{Lemma}

\theoremstyle{definition}
\newtheorem{definition}[subsection]{Definition}
\newtheorem{example}[subsection]{Example}
\newtheorem{exercise}[subsection]{Exercise}
\newtheorem{situation}[subsection]{Situation}

\theoremstyle{remark}
\newtheorem{remark}[subsection]{Remark}
\newtheorem{remarks}[subsection]{Remarks}

\numberwithin{equation}{subsection}

% Macros
%
\def\lim{\mathop{\mathrm{lim}}\nolimits}
\def\colim{\mathop{\mathrm{colim}}\nolimits}
\def\Spec{\mathop{\mathrm{Spec}}}
\def\Hom{\mathop{\mathrm{Hom}}\nolimits}
\def\Ext{\mathop{\mathrm{Ext}}\nolimits}
\def\SheafHom{\mathop{\mathcal{H}\!\mathit{om}}\nolimits}
\def\SheafExt{\mathop{\mathcal{E}\!\mathit{xt}}\nolimits}
\def\Sch{\mathit{Sch}}
\def\Mor{\mathop{\mathrm{Mor}}\nolimits}
\def\Ob{\mathop{\mathrm{Ob}}\nolimits}
\def\Sh{\mathop{\mathit{Sh}}\nolimits}
\def\NL{\mathop{N\!L}\nolimits}
\def\CH{\mathop{\mathrm{CH}}\nolimits}
\def\proetale{{pro\text{-}\acute{e}tale}}
\def\etale{{\acute{e}tale}}
\def\QCoh{\mathit{QCoh}}
\def\Ker{\mathop{\mathrm{Ker}}}
\def\Im{\mathop{\mathrm{Im}}}
\def\Coker{\mathop{\mathrm{Coker}}}
\def\Coim{\mathop{\mathrm{Coim}}}

% Boxtimes
%
\DeclareMathSymbol{\boxtimes}{\mathbin}{AMSa}{"02}

%
% Macros for moduli stacks/spaces
%
\def\QCohstack{\mathcal{QC}\!\mathit{oh}}
\def\Cohstack{\mathcal{C}\!\mathit{oh}}
\def\Spacesstack{\mathcal{S}\!\mathit{paces}}
\def\Quotfunctor{\mathrm{Quot}}
\def\Hilbfunctor{\mathrm{Hilb}}
\def\Curvesstack{\mathcal{C}\!\mathit{urves}}
\def\Polarizedstack{\mathcal{P}\!\mathit{olarized}}
\def\Complexesstack{\mathcal{C}\!\mathit{omplexes}}
% \Pic is the operator that assigns to X its picard group, usage \Pic(X)
% \Picardstack_{X/B} denotes the Picard stack of X over B
% \Picardfunctor_{X/B} denotes the Picard functor of X over B
\def\Pic{\mathop{\mathrm{Pic}}\nolimits}
\def\Picardstack{\mathcal{P}\!\mathit{ic}}
\def\Picardfunctor{\mathrm{Pic}}
\def\Deformationcategory{\mathcal{D}\!\mathit{ef}}
\usepackage{cleveref}
\crefname{originalsection}{original section}{original sections}
\Crefname{originalsection}{Original section}{Original sections}
\makeatletter
\newlabel{properties-example-does-not-work-in-general}{{01PN}{source}{}{sourcepointer-1}{}}
\newlabel{properties-example-does-not-work-in-general@cref}{{[originalsection][0][]01PN}{[1][0][]source}{}{}{}}
\makeatother
\crefname{originallemma}{original lemma}{original lemmas}
\Crefname{originallemma}{Original lemma}{Original lemmas}
\makeatletter
\newlabel{properties-lemma-sections-over-quasi-compact-open-in-affine}{{01PM}{source}{}{sourcepointer-2}{}}
\newlabel{properties-lemma-sections-over-quasi-compact-open-in-affine@cref}{{[originallemma][0][]01PM}{[1][0][]source}{}{}{}}
\makeatother
\title{An injective module whose associated sheaf on an affine scheme is not flasque}
\author{The Stacks Project Authors}
\date{stacks-project:examples.tex:local-snapshot-20261012; GNU FDL 1.2 or later}
\begin{document}
\maketitle
\noindent\textit{Editorial note.} All delivered mathematical text is the author's own, copied from the pinned Stacks Project chapter file examples.tex; no statement or proof step was rewritten, reordered or removed.\par\medskip
\noindent\textit{Editorial note.} The counted proof is the author's complete argument: The counted proof is 'See above'; the final standard-open restriction failure is the sharper second sentence of the statement, witnessed by U = D(x).\par\medskip
\noindent\textit{Editorial note.} The label lemma-nonflasque names the counted statement, and the proof environment immediately following it in the source is delivered verbatim as its proof.\par\medskip
\noindent\textit{Editorial note.} The counted statement is carried by the author's own lemma environment in the chapter source.\par\medskip
\noindent\textit{Allowed standard prerequisites (editorial).}\begin{itemize}
\item Context span 3734-3743: X = Spec(A) with A = k[x, y, z\_1, z\_2, ...]/(x\textasciicircum{}n z\_n), the ring from Properties, Example does-not-work-in-general; I = (x); U = D(x) quasi-compact open; as seen in the referenced Properties example there is a section s of O\_X(U) which does not come from an A-module map I\textasciicircum{}n -> A for any n >= 0
\item Context span 3745-3766: alpha: A -> J an embedding into an injective A-module, Q = J/alpha(A) with quotient beta; the claim is that Gamma(X, J\textasciitilde{}) -> Gamma(U, J\textasciitilde{}) is not surjective since alpha(s) is not in the image; by contradiction, x in J restricting to alpha(s) over U gives beta(x) in Q restricting to 0 over U, hence I\textasciicircum{}n beta(x) = 0 for some n (Properties, sections-over-quasi-compact-open-in-affine); this yields phi: I\textasciicircum{}n -> A, h -> alpha\textasciicircum{}\{-1\}(h x), which by the same Properties lemma gives rise to the section s - contradicting the choice of s
\end{itemize}\medskip
\ifdefined\hypertarget\hypertarget{sourcepointer-1}{}\fi
\noindent\textit{Source reference.} Original section 01PN: sources/stacks-chow/tags.txt, source lines 2236--2236. Prerequisite pointer invoked at 3784; it names a standard fact of a sibling chapter that the counted argument uses, and every step of the counted proof built on it is retained verbatim in the delivered body.\par\medskip
\ifdefined\hypertarget\hypertarget{sourcepointer-2}{}\fi
\noindent\textit{Source reference.} Original lemma 01PM: sources/stacks-chow/tags.txt, source lines 2235--2235. Prerequisite pointer invoked at 3805; it names a standard fact of a sibling chapter that the counted argument uses, and every step of the counted proof built on it is retained verbatim in the delivered body.\par\medskip
Consider the affine scheme $X = \Spec(A)$ where
$$
A = k[x, y, z_1, z_2, \ldots]/(x^nz_n)
$$
is the ring from
Properties, Example \ref{properties-example-does-not-work-in-general}.
Set $I = (x) \subset A$. Consider the quasi-compact open $U = D(x)$ of $X$.
We have seen in loc.\ cit.\ that there is a section
$s \in \mathcal{O}_X(U)$ which does not come from an $A$-module
map $I^n \to A$ for any $n \geq 0$.

\medskip\noindent
Let $\alpha : A \to J$ be the embedding of $A$ into an injective $A$-module.
Let $Q = J/\alpha(A)$ and denote $\beta : J \to Q$ the quotient map.
We claim that the map
$$
\Gamma(X, \widetilde{J})
\longrightarrow
\Gamma(U, \widetilde{J})
$$
is not surjective. Namely, we claim that $\alpha(s)$ is not in the image.
To see this, we argue by contradiction. So assume that $x \in J$ is an
element which restricts to $\alpha(s)$ over $U$. Then $\beta(x) \in Q$
is an element which restricts to $0$ over $U$. Hence we know that
$I^n\beta(x) = 0$ for some $n$, see
Properties,
Lemma \ref{properties-lemma-sections-over-quasi-compact-open-in-affine}.
This implies that we get a morphism
$\varphi : I^n \to A$, $h \mapsto \alpha^{-1}(hx)$. It is easy to see that
this morphism $\varphi$ gives rise to the section $s$ via the map of
Properties,
Lemma \ref{properties-lemma-sections-over-quasi-compact-open-in-affine}
which is a contradiction.

\begin{lemma}
\label{lemma-nonflasque}
There exists an affine scheme $X = \Spec(A)$ and an injective
$A$-module $J$ such that $\widetilde{J}$ is not a flasque sheaf on $X$.
Even the restriction $\Gamma(X, \widetilde{J}) \to \Gamma(U, \widetilde{J})$
with $U$ a standard open need not be surjective.
\end{lemma}

\begin{proof}
See above.
\end{proof}

\end{document}
